% Generated draft: requires mathematical and source review.
\documentclass[11pt]{article}
\usepackage[margin=25mm]{geometry}
\usepackage{fontspec}
\setmainfont{Latin Modern Roman}
\usepackage{amsmath,amssymb,mathtools,unicode-math}
\setmathfont{Latin Modern Math}
\usepackage{enumitem,needspace,xurl,hyperref}
\setlength{\parindent}{0pt}
\setlength{\parskip}{5pt}
\setlength{\emergencystretch}{3em}
\newcommand{\N}{\mathbb{N}}
\newcommand{\Z}{\mathbb{Z}}
\newcommand{\Q}{\mathbb{Q}}
\newcommand{\R}{\mathbb{R}}
\newcommand{\C}{\mathbb{C}}
\newcommand{\F}{\mathbb{F}}
\newcommand{\A}{\mathbb{A}}
\newcommand{\PP}{\mathbb{P}}
\newcommand{\E}{\mathbb{E}}
\newcommand{\eps}{\varepsilon}
\newcommand{\dd}{\,\mathrm{d}}
\begin{document}
\section*{UnsolvedMath ID: 30004713}
\textbf{Source problem code:} catalogue entry
\subsection{Definitions and mathematical statement}
Under the standard interpretation in terms of holomorphic Abelian differentials with all zeros marked, let \(g\geq 2\), let \(n\geq 1\), and let \(\mu=(m_1,\ldots,m_n)\) be a tuple of positive integers satisfying \(\sum_{i=1}^n m_i=2g-2\). Define the stratum
\[
\mathcal H_g(\mu)=\left\{(C,p_1,\ldots,p_n,\omega):\begin{array}{l}C\text{ is a smooth connected projective complex curve of genus }g,\\ p_1,\ldots,p_n\text{ are distinct ordered points},\\0\neq\omega\in H^0(C,K_C),\quad \operatorname{div}(\omega)=\sum_{i=1}^n m_i p_i\end{array}\right\}/\cong.
\]
Here \(K_C\) is the canonical line bundle, and isomorphisms preserve the marked points and the differential. Its projectivization \(\mathbb P\mathcal H_g(\mu)\) identifies \(\omega\) with \(a\omega\) for every \(a\in\C^*\). These definitions make explicit a conventional interpretation; the supplied record does not specify the ground field, whether every marking must be a zero, or whether markings are ordered.

For an irreducible complex algebraic variety \(X\), its ordinary birational Kodaira dimension means \(\kappa(\widetilde X,K_{\widetilde X})\), where \(\widetilde X\) is a smooth projective model of its function field. This invariant is \(-\infty\) if all positive pluricanonical systems are empty, and otherwise is the maximum dimension of the images of the pluricanonical rational maps. General type means that this dimension equals \(\dim X\). For a nonproper variety, logarithmic Kodaira dimension instead uses \(K_{\overline X}+D\) on a smooth projective compactification with simple normal crossings boundary \(D\); it need not equal the ordinary invariant.

\textbf{Problem as supplied, with unresolved choices retained.} Determine the behavior, as \(g\) becomes large, of the Kodaira dimensions of moduli spaces of Abelian differentials, especially strata having few marked points. Once the intended moduli space and Kodaira-dimension convention are fixed, this asks for the relevant dimensions on its irreducible components as functions of \(g\) and the stratum data \(\mu\). The record does not specify whether the target is \(\mathcal H_g(\mu)\), its projectivization, a locus of pointed curves, or a particular compactification, nor does it specify a quantitative meaning of ``few'' or an asymptotic quantifier. Consequently, no assertion of eventual general type for every such stratum is supplied.

The explicitly distinguished case is \(n=1\) and \(\mu=(2g-2)\). In the interpretation above, it consists of differentials with one zero of order \(2g-2\). After identifying scalar multiples of the differential, its points correspond to subcanonical pointed curves
\[
\mathcal S_g=\{(C,p): C\text{ is a smooth connected projective complex curve of genus }g,\ K_C\cong\mathcal O_C((2g-2)p)\}/\cong.
\]
The question includes the large-genus Kodaira-dimension behavior in this case. The intended components, treatment of hyperelliptic loci, and precise target invariant remain unspecified.
\subsection{Short English statement}
Determine how the Kodaira dimensions of strata of Abelian differentials behave in large genus, particularly when there are few marked zeros, including the one-zero stratum \(\mu=(2g-2)\). A precise componentwise question requires specifying the moduli space, the Kodaira-dimension convention, and what ``few'' means. The supplied partially solved label and literature assessment are dated curation requiring independent review.
\subsection{Sources}
\begin{itemize}
\item Local UnsolvedMath registry: supplied record for permanent ID \(30004713\), title ``Kodaira Dimensions of Abelian-Differential Moduli Spaces,'' exact source problem code \texttt{catalogue entry}, in the Oberwolfach Reports Open Problems collection. The record gives proposed year \(2021\) and names no proposer.
\item Supplied source citation: \emph{Classical Algebraic Geometry (2022)}, Oberwolfach Reports, DOI \texttt{10.4171/owr/2021/29}, pp.\ 1519--1577. \url{https://doi.org/10.4171/owr/2021/29}. The background describes Oberwolfach Reports as documenting mathematical workshops and their open-problem discussions. The citation's year \(2022\), the proposed year \(2021\), and the DOI containing \texttt{2021} are retained as supplied.
\item Supplied literature source: Chen--Costantini--Möller, \url{https://arxiv.org/abs/2204.11943}.
\item Supplied literature source: Bud--Chen--Möller, \url{https://arxiv.org/abs/2410.18719}.
\item Literature assessment checked at \texttt{2026-08-22}: the supplied status is ``partially solved.'' The supplied assessment states that general type is proved for many few-zero strata and for the large-genus odd and even minimal spin components, while the broad few-point Kodaira-dimension question remains unresolved. This is supplied dated curation derived from the OWR record, not an independently verified literature assessment; the time-sensitive status and cited claims require independent verification.
\item The supplied statement status is \texttt{corrected\_verified}. Its accompanying verification note says that the genus, marking, and stratum parameters were normalized as LaTeX without changing the problem. This records the supplied editorial assessment and does not establish independent mathematical verification of this draft.
\end{itemize}
\end{document}
